% =============================================================
% Capitulo 1 - Introduccion
% =============================================================
\chapter{Introducci\'on}
\label{cap:introduccion}

\section{Antecedentes del problema}
\label{sec:antecedentes}

En 2013 se dio inicio al proyecto transnacional entre Argentina, Chile
y Jap\'on para el monitoreo y estudio de la atm\'osfera en el cono sur
de Sudam\'erica, en cuyo marco se establecieron siete nodos de
observaci\'on en el territorio argentino, dotados de instrumental
espec\'ifico para el monitoreo de par\'ametros atmosf\'ericos
\cite{sat2026}. En 2017 se puso en marcha el nodo de observaci\'on del
Observatorio Geof\'isico y Meteorol\'ogico de Pilar, C\'ordoba, del
Servicio Meteorol\'ogico Nacional (SMN), sitio caracterizado por las
frecuentes irrupciones de aerosoles provenientes de la quema de biomasa
en el norte argentino, Bolivia, Paraguay y el sur de Brasil
\cite{sat2026}. Dada esta particularidad geogr\'afica, el nodo fue
dotado de un LiDAR de Alta Resoluci\'on Espectral (HSRL, por sus siglas
en ingl\'es), puesto en funcionamiento durante la pr\'actica
profesional supervisada (PPS) realizada por el autor en 2023
\cite{pps2025}.

Si bien el instrumento alcanz\'o un funcionamiento \'optimo al
finalizar la PPS, la naturaleza experimental de su dise\~no gener\'o
inestabilidades para la operaci\'on a largo plazo: el software
principal no posee capacidad de resoluci\'on de fallos, la PC en la que
corre el sistema tiene capacidades de c\'omputo limitadas, la
implementaci\'on actual disminuye la vida \'util de la l\'ampara flash
del l\'aser debido al elevado n\'umero de disparos requeridos, y el
sistema depende de un costoso osciloscopio Tektronix de cuatro canales
con software desactualizado \cite{sat2026}.

% PLACEHOLDER: ampliar con el detalle especifico de fallas u operacion
% reportadas durante el uso del sistema heredado de la PPS, si se cuenta
% con registros adicionales.

\section{Relevancia del trabajo}
\label{sec:relevancia}

El HSRL de Pilar constituye uno de los instrumentos de mayor complejidad
de la red SAVER-Net, integrada por ocho nodos de observaci\'on en
Argentina y uno en Chile \cite{galvagno2026ilrc}, cuya operaci\'on
responde a compromisos internacionales asumidos en el marco del
programa SATREPS, financiado por la Agencia de Cooperaci\'on
Internacional del Jap\'on (JICA) y la Agencia de Ciencia y Tecnolog\'ia
de Jap\'on (JST) \cite{galvagno2026ilrc}. La informaci\'on generada por
el canal de alta resoluci\'on espectral resulta relevante a nivel
nacional e internacional para el estudio de aerosoles atmosf\'ericos, lo
que fue reconocido en la Solicitud de Aprobaci\'on de Tema del presente
proyecto como de impacto ambiental directo \cite{sat2026}.

Avances preliminares de este trabajo fueron presentados en la 32.\textsuperscript{a}
Conferencia Internacional de Radar L\'aser (ILRC32) \cite{galvagno2026ilrc},
lo que da cuenta del inter\'es de la comunidad cient\'ifica internacional
en la automatizaci\'on del lazo de sintonizaci\'on del HSRL.

\section{Motivaci\'on para la elecci\'on del tema}
\label{sec:motivacion}

La motivaci\'on del proyecto surge directamente de las limitaciones
identificadas en el sistema de control utilizado durante la PPS
\cite{pps2025}. La dependencia de un osciloscopio y una PC de
escritorio para realizar el lazo de control de sintonizaci\'on
introduce un punto de falla poco robusto, encarece el mantenimiento del
sistema y limita la posibilidad de operaci\'on remota y desatendida del
instrumento. Reemplazar este conjunto por una placa electr\'onica
dedicada, basada en un microcontrolador, permitir\'ia dotar al sistema
de mayor robustez y funcionalidad \cite{sat2026}.

% PLACEHOLDER: agregar motivacion personal/profesional del autor
% vinculada a su rol como responsable tecnico del nodo SaverNet-Cordoba.

\section{Formulaci\'on del problema}
\label{sec:formulacion_problema}

El problema a resolver consiste en dise\~nar y construir un sistema
electr\'onico capaz de reemplazar al conjunto osciloscopio-PC-software
actualmente utilizado en el lazo de control de sintonizaci\'on del
canal HSR, de modo de automatizar la medici\'on de las se\~nales
provenientes de los fotodetectores y la toma de decisiones para
modificar la longitud de onda del l\'aser seeder \cite{sat2026}. El
sistema debe acondicionar y digitalizar se\~nales pulsadas de corta
duraci\'on (del orden de los nanosegundos), ejecutar un algoritmo de
control que permita igualar las se\~nales $P_+$ y $P_-$ del lazo,
comunicarse con el l\'aser seeder mediante el protocolo definido por su
fabricante, e informar su estado a un operador a trav\'es de una
interfaz de usuario \cite{sat2026}.

\section{Objetivo general}
\label{sec:objetivo_general}

Desarrollar y ensamblar un sistema electr\'onico comandado por PC que
permita controlar y automatizar los procesos de sintonizaci\'on de
longitud de onda del LiDAR de alta resoluci\'on espectral (HSRL)
presente en el Observatorio Geof\'isico y Meteorol\'ogico Pilar del
Servicio Meteorol\'ogico Nacional, en reemplazo del conjunto compuesto
por osciloscopio, software y PC actualmente utilizado, brindando mayor
robustez y funcionalidad al sistema \cite{sat2026}.

\section{Objetivos espec\'ificos}
\label{sec:objetivos_especificos}

De acuerdo a lo establecido en la Solicitud de Aprobaci\'on de Tema
\cite{sat2026}, los objetivos espec\'ificos del proyecto son:

\begin{itemize}
  \item Dise\~nar un circuito capaz de detectar se\~nales pulsadas de
        alta velocidad y acondicionarlas para su posterior
        digitalizaci\'on.
  \item Implementar y validar el circuito de detecci\'on de acuerdo a
        las se\~nales medidas por el bucle de control.
  \item Seleccionar los conversores anal\'ogico/digital y el
        microcontrolador adecuados para la digitalizaci\'on y el
        procesamiento de las se\~nales pulsadas.
  \item Programar el microcontrolador para efectuar la digitalizaci\'on
        de la se\~nal previamente acondicionada mediante los
        conversores seleccionados.
  \item Implementar y validar la comunicaci\'on entre el l\'aser seeder
        y el microcontrolador del sistema para efectuar la selecci\'on
        de longitud de onda.
  \item Desarrollar un software para el microcontrolador que tome
        decisiones a partir de los datos medidos.
  \item Programar una interfaz m\'aquina-humano (HMI) en Python para
        permitir la interacci\'on entre el PC y el microcontrolador.
  \item Validar el correcto funcionamiento del sistema de alta
        resoluci\'on espectral (HSR) a partir de las mediciones
        obtenidas con el sistema LiDAR.
\end{itemize}

\section{Metodolog\'ia utilizada para lograr los objetivos propuestos}
\label{sec:metodologia_intro}

El desarrollo del proyecto se organiz\'o en las siguientes etapas
\cite{sat2026}:

\begin{enumerate}
  \item Investigaci\'on y revisi\'on bibliogr\'afica de los conceptos
        te\'oricos involucrados.
  \item Dise\~no, implementaci\'on y validaci\'on del circuito de
        detecci\'on de se\~nales.
  \item Selecci\'on y ensamble del conjunto ADC/microcontrolador.
  \item Desarrollo del firmware del microcontrolador.
  \item Desarrollo del software HMI para la comunicaci\'on entre la PC
        y el microcontrolador.
  \item Validaci\'on del funcionamiento del sistema completo a partir
        de mediciones reales sobre el instrumento LiDAR.
\end{enumerate}

Al cierre del presente informe, las etapas 1 a 3 se encuentran
finalizadas y las etapas 4 y 5 se encuentran en desarrollo, restando la
etapa de validaci\'on integral del sistema (ver
Cap\'itulo~\ref{cap:modelo_experimental} y apartado de
Resultados)\footnote{Estado de avance seg\'un los Informes de Avance
Mensual correspondientes a enero-febrero, marzo y abril de 2026
\cite{informeEneFeb2026,informeMarzo2026,informeAbril2026}.}.

\section{Organizaci\'on del informe}
\label{sec:organizacion_informe}

El presente informe se organiza en tres partes. La Parte~I (Marco
Te\'orico, Cap\'itulos~\ref{cap:lidar_hsrl} a \ref{cap:control_hmi})
desarrolla los fundamentos te\'oricos de la tecnolog\'ia LiDAR, del
HSRL, de la fotodetecci\'on de se\~nales pulsadas, de los sistemas
embebidos utilizados y del control autom\'atico aplicado al problema
de sintonizaci\'on. La Parte~II (Marco Metodol\'ogico,
Cap\'itulos~\ref{cap:diseno_deteccion_picos} a
\ref{cap:modelo_experimental}) describe el dise\~no, la
implementaci\'on y la validaci\'on del circuito de detecci\'on, de la
placa de control, del firmware del microcontrolador y del software de
interfaz, as\'i como el modelo experimental resultante. Finalmente, la
Parte~III presenta los resultados obtenidos y las conclusiones del
trabajo. Los Anexos re\'unen las hojas de datos de componentes, el
manual de comunicaci\'on del seeder, el c\'odigo fuente relevante y la
documentaci\'on presentada ante la C\'atedra Proyecto Integrador
(Solicitud de Aprobaci\'on de Tema e Informes de Avance Mensual).
